% !TeX program = xelatex
% 编译时以 book/ 为工作目录；正文按主题拆分在 content/。
\documentclass[UTF8,zihao=-4,oneside,openany,fontset=fandol]{ctexbook}

\usepackage[a4paper,margin=2.6cm,headheight=15pt]{geometry}
\usepackage{amsmath,amssymb}
\usepackage{graphicx}
\usepackage{booktabs}
\usepackage[unicode,hidelinks]{hyperref}

\graphicspath{{../figures/}}
\setcounter{secnumdepth}{2}
\setcounter{tocdepth}{2}
% 空白段落让仅含标题的框架正常分页；填写正文时删除对应占位命令。
\newcommand{\contentplaceholder}{\par\noindent\mbox{}\par}
\hypersetup{
  pdftitle={生成式人工智能：从概率建模到扩散模型},
  bookmarksnumbered=true,
  bookmarksdepth=2
}

\title{生成式人工智能\\[0.5em]\large 从概率建模到扩散模型}
\author{}
\date{}

\begin{document}
\raggedbottom

\frontmatter
\hypersetup{pageanchor=false}
\maketitle
\hypersetup{pageanchor=true}
\tableofcontents
% 前置部分保留目录条目，但不出现 0.1 等正文式编号。
\setcounter{secnumdepth}{-1}
% 来源：README.md
% 仅建立章节框架；正文、公式、图表与习题内容待后续迁移。
\chapter{阅读指南}
\label{chap:reading-guide}

\section{从哪里开始？}
\label{sec:reading-guide-routes}
% TODO：补充本节内容。
\contentplaceholder

\section{知识目录与章节关系}
\label{sec:reading-guide-roadmap}
% TODO：补充本节内容。
\contentplaceholder

\subsection{从概率建模到五类生成模型}
\label{subsec:reading-guide-connections}
% TODO：补充本节内容。
\contentplaceholder

\subsection{模型表示、学习目标与采样成本}
\label{subsec:reading-guide-comparison}
% TODO：补充本节内容。
\contentplaceholder

\section{阅读约定}
\label{sec:reading-guide-conventions}
% TODO：补充本节内容。
\contentplaceholder

% 来源：docs/notation.md
% 仅建立章节框架；正文、公式、图表与习题内容待后续迁移。
\chapter{符号约定与数学阅读提示}
\label{chap:notation}

\section{数据、模型与目标}
\label{sec:notation-data}
% TODO：补充本节内容。
\contentplaceholder

\section{概率与信息量}
\label{sec:notation-probability}
% TODO：补充本节内容。
\contentplaceholder

\subsection{联合分布、条件分布与边缘分布}
\label{subsec:notation-distributions}
% TODO：补充本节内容。
\contentplaceholder

\subsection{熵、交叉熵与分布差异}
\label{subsec:notation-divergences}
% TODO：补充本节内容。
\contentplaceholder

\subsection{概率质量与概率密度}
\label{subsec:notation-mass-density}
% TODO：补充本节内容。
\contentplaceholder

\section{高斯与线性代数}
\label{sec:notation-linear-algebra}
% TODO：补充本节内容。
\contentplaceholder

\subsection{高斯分布的参数约定}
\label{subsec:notation-gaussian}
% TODO：补充本节内容。
\contentplaceholder

\subsection{向量、矩阵与微分算子}
\label{subsec:notation-operators}
% TODO：补充本节内容。
\contentplaceholder

\subsection{对数方差与对数标准差}
\label{subsec:notation-log-scales}
% TODO：补充本节内容。
\contentplaceholder

\section{各模型的专用约定}
\label{sec:notation-model-conventions}
% TODO：补充本节内容。
\contentplaceholder

\section{读公式时先问四个问题}
\label{sec:notation-reading-formulas}
% TODO：补充本节内容。
\contentplaceholder


\mainmatter
\setcounter{secnumdepth}{2}
% 来源：docs/01-introduction.md
% 仅建立章节框架；正文、公式、图表与习题内容待后续迁移。
\chapter{生成建模在解决什么问题？}
\label{chap:intro}

\section*{本章要点}
% TODO：迁移原章“先抓住这几件事”的要点。
\contentplaceholder

\section{用体重生成理解整个流程}
\label{sec:intro-workflow}
% TODO：补充本节内容。
\contentplaceholder

\subsection{数据分布与模型族}
\label{subsec:intro-distribution}
% TODO：补充本节内容。
\contentplaceholder

\subsection{参数学习、采样与检查}
\label{subsec:intro-learning-sampling}
% TODO：补充本节内容。
\contentplaceholder

\section{学到分布以后能做什么？}
\label{sec:intro-applications}
% TODO：补充本节内容。
\contentplaceholder

\subsection{采样与条件生成}
\label{subsec:intro-sampling}
% TODO：补充本节内容。
\contentplaceholder

\subsection{似然估计及其边界}
\label{subsec:intro-likelihood}
% TODO：补充本节内容。
\contentplaceholder

\subsection{潜变量与表征学习}
\label{subsec:intro-representation}
% TODO：补充本节内容。
\contentplaceholder

\section{与分类、图形学有什么关系？}
\label{sec:intro-related-tasks}
% TODO：补充本节内容。
\contentplaceholder

\subsection{无条件生成、条件生成与分类}
\label{subsec:intro-classification}
% TODO：补充本节内容。
\contentplaceholder

\subsection{统计生成模型与传统图形学}
\label{subsec:intro-graphics}
% TODO：补充本节内容。
\contentplaceholder

\section{为什么高维生成困难？}
\label{sec:intro-challenges}
% TODO：补充本节内容。
\contentplaceholder

\subsection{高维表征与复杂分布建模}
\label{subsec:intro-representation-modeling}
% TODO：补充本节内容。
\contentplaceholder

\subsection{学习准则与评价需求}
\label{subsec:intro-learning-evaluation}
% TODO：补充本节内容。
\contentplaceholder

\section{课程范围与阅读准备}
\label{sec:intro-scope}
% TODO：补充本节内容。
\contentplaceholder

\subsection{先修知识与 AI for Science 示例}
\label{subsec:intro-prerequisites}
% TODO：补充本节内容。
\contentplaceholder

\subsection{生成建模的发展线索}
\label{subsec:intro-history}
% TODO：补充本节内容。
\contentplaceholder

\section{自检}
\label{sec:intro-review}
% TODO：补充本节内容。
\contentplaceholder

\section{来源与延伸阅读}
\label{sec:intro-reading}
% TODO：补充本节内容。
\contentplaceholder

% 来源：docs/02-probabilistic-modeling.md
% 仅建立章节框架；正文、公式、图表与习题内容待后续迁移。
\chapter{概率建模：从联合分布到结构化表示}
\label{chap:prob}

\section*{本章要点}
% TODO：迁移原章“先抓住这几件事”的要点。
\contentplaceholder

\section{概率学习的三个组成}
\label{sec:prob-components}
% TODO：补充本节内容。
\contentplaceholder

\subsection{模型族、学习准则与优化算法}
\label{subsec:prob-family}
% TODO：补充本节内容。
\contentplaceholder

\subsection{总体目标与样本近似}
\label{subsec:prob-population-empirical}
% TODO：补充本节内容。
\contentplaceholder

\section{直接表示联合分布的代价}
\label{sec:prob-joint-cost}
% TODO：补充本节内容。
\contentplaceholder

\subsection{概率归一化与自由参数}
\label{subsec:prob-parameters}
% TODO：补充本节内容。
\contentplaceholder

\subsection{联合状态空间与维度灾难}
\label{subsec:prob-dimensionality}
% TODO：补充本节内容。
\contentplaceholder

\section{三种分解，三种表达能力}
\label{sec:prob-factorizations}
% TODO：补充本节内容。
\contentplaceholder

\subsection{完整链式分解}
\label{subsec:prob-chain-rule}
% TODO：补充本节内容。
\contentplaceholder

\subsection{完全独立}
\label{subsec:prob-independence}
% TODO：补充本节内容。
\contentplaceholder

\subsection{条件独立与一阶马尔可夫链}
\label{subsec:prob-markov}
% TODO：补充本节内容。
\contentplaceholder

\section{贝叶斯网络把假设画成图}
\label{sec:prob-bayesian-networks}
% TODO：补充本节内容。
\contentplaceholder

\subsection{有向无环图与局部分解}
\label{subsec:prob-dag}
% TODO：补充本节内容。
\contentplaceholder

\subsection{条件概率表与参数计数}
\label{subsec:prob-cpt}
% TODO：补充本节内容。
\contentplaceholder

\subsection{例子：难度、智力、成绩与推荐信}
\label{subsec:prob-student-example}
% TODO：补充本节内容。
\contentplaceholder

\subsection{拓扑顺序与祖先采样}
\label{subsec:prob-ancestral-sampling}
% TODO：补充本节内容。
\contentplaceholder

\section{生成式分类与判别式分类}
\label{sec:prob-classification}
% TODO：补充本节内容。
\contentplaceholder

\subsection{朴素贝叶斯：建模联合，再推导后验}
\label{subsec:prob-naive-bayes}
% TODO：补充本节内容。
\contentplaceholder

\subsection{逻辑回归：直接建模类别后验}
\label{subsec:prob-logistic-regression}
% TODO：补充本节内容。
\contentplaceholder

\subsection{建模对象与决策边界}
\label{subsec:prob-boundaries}
% TODO：补充本节内容。
\contentplaceholder

\section{分解结构与参数化方式要分开看}
\label{sec:prob-parameterization}
% TODO：补充本节内容。
\contentplaceholder

\subsection{条件概率表与神经网络参数化}
\label{subsec:prob-conditional-models}
% TODO：补充本节内容。
\contentplaceholder

\subsection{条件集合、参数共享与表达能力}
\label{subsec:prob-complexity}
% TODO：补充本节内容。
\contentplaceholder

\section{自检}
\label{sec:prob-review}
% TODO：补充本节内容。
\contentplaceholder

\section{来源与延伸阅读}
\label{sec:prob-reading}
% TODO：补充本节内容。
\contentplaceholder

% 来源：docs/03-autoregressive-models.md
% 仅建立章节框架；正文、公式、图表与习题内容待后续迁移。
\chapter{自回归模型}
\label{chap:ar}

\section*{本章要点}
% TODO：迁移原章“先抓住这几件事”的要点。
\contentplaceholder

\section{先确定顺序，再建模条件分布}
\label{sec:ar-ordering}
% TODO：补充本节内容。
\contentplaceholder

\subsection{联合概率分解与归一化}
\label{subsec:ar-chain-rule}
% TODO：补充本节内容。
\contentplaceholder

\subsection{变量顺序与归纳偏置}
\label{subsec:ar-inductive-bias}
% TODO：补充本节内容。
\contentplaceholder

\section{网络输出的是分布参数}
\label{sec:ar-outputs}
% TODO：补充本节内容。
\contentplaceholder

\subsection{二元变量与伯努利输出}
\label{subsec:ar-binary}
% TODO：补充本节内容。
\contentplaceholder

\subsection{类别变量与连续变量的输出}
\label{subsec:ar-categorical-continuous}
% TODO：补充本节内容。
\contentplaceholder

\subsection{概率参数、采样与确定性决策}
\label{subsec:ar-sampling-decisions}
% TODO：补充本节内容。
\contentplaceholder

\section{从简单模型到共享网络}
\label{sec:ar-shared-networks}
% TODO：补充本节内容。
\contentplaceholder

\subsection{FVSBN：用逻辑回归表示条件概率}
\label{subsec:ar-fvsbn}
% TODO：补充本节内容。
\contentplaceholder

\subsection{NADE：非线性条件分布与参数共享}
\label{subsec:ar-nade}
% TODO：补充本节内容。
\contentplaceholder

\subsection{MADE：在一个网络内施加掩码}
\label{subsec:ar-made}
% TODO：补充本节内容。
\contentplaceholder

\section{架构不同，因果约束相同}
\label{sec:ar-architectures}
% TODO：补充本节内容。
\contentplaceholder

\subsection{RNN 与 PixelRNN}
\label{subsec:ar-rnn}
% TODO：补充本节内容。
\contentplaceholder

\subsection{因果 CNN 与 WaveNet}
\label{subsec:ar-causal-cnn}
% TODO：补充本节内容。
\contentplaceholder

\subsection{PixelCNN 的空间与通道掩码}
\label{subsec:ar-pixelcnn}
% TODO：补充本节内容。
\contentplaceholder

\subsection{Transformer 的输入移位与因果注意力}
\label{subsec:ar-transformer}
% TODO：补充本节内容。
\contentplaceholder

\section{生成、似然计算与训练要分开看}
\label{sec:ar-computation}
% TODO：补充本节内容。
\contentplaceholder

\subsection{生成：祖先采样}
\label{subsec:ar-generation}
% TODO：补充本节内容。
\contentplaceholder

\subsection{似然计算：观测已经完整}
\label{subsec:ar-likelihood}
% TODO：补充本节内容。
\contentplaceholder

\subsection{训练：Teacher Forcing}
\label{subsec:ar-teacher-forcing}
% TODO：补充本节内容。
\contentplaceholder

\subsection{并行计算、曝光偏差与误差积累}
\label{subsec:ar-parallelism}
% TODO：补充本节内容。
\contentplaceholder

\section{这类模型适合什么？}
\label{sec:ar-strengths-limits}
% TODO：补充本节内容。
\contentplaceholder

\subsection{似然建模的优势与顺序生成的代价}
\label{subsec:ar-tradeoffs}
% TODO：补充本节内容。
\contentplaceholder

\subsection{连续数据与潜在 token 表示}
\label{subsec:ar-representations}
% TODO：补充本节内容。
\contentplaceholder

\section{自检}
\label{sec:ar-review}
% TODO：补充本节内容。
\contentplaceholder

\section{来源与延伸阅读}
\label{sec:ar-reading}
% TODO：补充本节内容。
\contentplaceholder

% 来源：docs/04-maximum-likelihood.md
% 仅建立章节框架；正文、公式、图表与习题内容待后续迁移。
\chapter{最大似然与泛化}
\label{chap:mle}

\section*{本章要点}
% TODO：迁移原章“先抓住这几件事”的要点。
\contentplaceholder

\section{用 KL 散度比较两个分布}
\label{sec:mle-kl}
% TODO：补充本节内容。
\contentplaceholder

\subsection{离散与连续情形的 KL 定义}
\label{subsec:mle-definition}
% TODO：补充本节内容。
\contentplaceholder

\subsection{支持集条件与有限性}
\label{subsec:mle-support}
% TODO：补充本节内容。
\contentplaceholder

\subsection{KL 非负性的推导与取等条件}
\label{subsec:mle-nonnegativity}
% TODO：补充本节内容。
\contentplaceholder

\subsection{熵、交叉熵与 KL 的关系}
\label{subsec:mle-entropy}
% TODO：补充本节内容。
\contentplaceholder

\section{从分布拟合得到最大似然}
\label{sec:mle-distribution-fitting}
% TODO：补充本节内容。
\contentplaceholder

\subsection{总体正向 KL 与期望对数似然}
\label{subsec:mle-population}
% TODO：补充本节内容。
\contentplaceholder

\subsection{正向与反向 KL 的模式拟合倾向}
\label{subsec:mle-direction}
% TODO：补充本节内容。
\contentplaceholder

\subsection{独立同分布样本与经验对数似然}
\label{subsec:mle-empirical}
% TODO：补充本节内容。
\contentplaceholder

\subsection{似然、对数似然与平均负对数似然}
\label{subsec:mle-nll}
% TODO：补充本节内容。
\contentplaceholder

\section{两个可手算的例子}
\label{sec:mle-examples}
% TODO：补充本节内容。
\contentplaceholder

\subsection{硬币：频率就是伯努利参数的 MLE}
\label{subsec:mle-bernoulli}
% TODO：补充本节内容。
\contentplaceholder

\subsection{高斯均值：已知方差下得到样本平均}
\label{subsec:mle-gaussian}
% TODO：补充本节内容。
\contentplaceholder

\section{复杂模型：分解目标，再用梯度优化}
\label{sec:mle-optimization}
% TODO：补充本节内容。
\contentplaceholder

\subsection{联合对数似然的局部分解}
\label{subsec:mle-local}
% TODO：补充本节内容。
\contentplaceholder

\subsection{离散条件概率表的局部估计}
\label{subsec:mle-cpt}
% TODO：补充本节内容。
\contentplaceholder

\subsection{梯度上升、梯度下降与非凸优化}
\label{subsec:mle-gradient}
% TODO：补充本节内容。
\contentplaceholder

\subsection{小批量梯度估计及其条件}
\label{subsec:mle-minibatch}
% TODO：补充本节内容。
\contentplaceholder

\section{样本平均为何有用，又为何会过拟合？}
\label{sec:mle-generalization}
% TODO：补充本节内容。
\contentplaceholder

\subsection{蒙特卡洛平均的期望与方差}
\label{subsec:mle-monte-carlo}
% TODO：补充本节内容。
\contentplaceholder

\subsection{大数定律与数据依赖的参数估计}
\label{subsec:mle-data-dependent}
% TODO：补充本节内容。
\contentplaceholder

\subsection{总体风险、经验风险与验证风险}
\label{subsec:mle-risk}
% TODO：补充本节内容。
\contentplaceholder

\subsection{欠拟合、过拟合与偏差—方差视角}
\label{subsec:mle-fitting}
% TODO：补充本节内容。
\contentplaceholder

\subsection{模型容量与泛化界的适用条件}
\label{subsec:mle-bounds}
% TODO：补充本节内容。
\contentplaceholder

\subsection{训练似然、验证似然与生成质量}
\label{subsec:mle-evaluation}
% TODO：补充本节内容。
\contentplaceholder

\section{自检}
\label{sec:mle-review}
% TODO：补充本节内容。
\contentplaceholder

\section{来源与延伸阅读}
\label{sec:mle-reading}
% TODO：补充本节内容。
\contentplaceholder

% 来源：docs/05-latent-variable-models-and-vae.md
% 仅建立章节框架；正文、公式、图表与习题内容待后续迁移。
\chapter{潜变量模型与变分自编码器}
\label{chap:vae}

\section*{本章要点}
% TODO：迁移原章“先抓住这几件事”的要点。
\contentplaceholder

\section{为什么引入看不见的变量？}
\label{sec:vae-motivation}
% TODO：补充本节内容。
\contentplaceholder

\subsection{联合分布、先验、条件分布与后验}
\label{subsec:vae-joint}
% TODO：补充本节内容。
\contentplaceholder

\subsection{高斯混合模型与分量后验}
\label{subsec:vae-gmm}
% TODO：补充本节内容。
\contentplaceholder

\subsection{连续潜变量与神经网络条件分布}
\label{subsec:vae-continuous}
% TODO：补充本节内容。
\contentplaceholder

\section{困难在哪里：边缘化与后验}
\label{sec:vae-inference}
% TODO：补充本节内容。
\contentplaceholder

\subsection{边缘似然与后验的计算困难}
\label{subsec:vae-marginalization}
% TODO：补充本节内容。
\contentplaceholder

\subsection{重要性采样与提议分布}
\label{subsec:vae-importance-sampling}
% TODO：补充本节内容。
\contentplaceholder

\subsection{似然估计与对数似然估计的区别}
\label{subsec:vae-log-bias}
% TODO：补充本节内容。
\contentplaceholder

\section{ELBO：将难算的对数似然变成可优化目标}
\label{sec:vae-elbo}
% TODO：补充本节内容。
\contentplaceholder

\subsection{从 Jensen 不等式推导证据下界}
\label{subsec:vae-jensen}
% TODO：补充本节内容。
\contentplaceholder

\subsection{后验 KL 间隙与变分推断}
\label{subsec:vae-posterior-gap}
% TODO：补充本节内容。
\contentplaceholder

\subsection{与 EM 算法的联系}
\label{subsec:vae-em}
% TODO：补充本节内容。
\contentplaceholder

\subsection{重建对数似然与先验正则}
\label{subsec:vae-decomposition}
% TODO：补充本节内容。
\contentplaceholder

\section{VAE：用两个网络实现变分学习}
\label{sec:vae-networks}
% TODO：补充本节内容。
\contentplaceholder

\subsection{编码器与摊销推断}
\label{subsec:vae-amortized}
% TODO：补充本节内容。
\contentplaceholder

\subsection{解码器与观测分布假设}
\label{subsec:vae-observation}
% TODO：补充本节内容。
\contentplaceholder

\subsection{对角高斯后验的闭式 KL}
\label{subsec:vae-gaussian-kl}
% TODO：补充本节内容。
\contentplaceholder

\subsection{负 ELBO 与训练损失}
\label{subsec:vae-negative-elbo}
% TODO：补充本节内容。
\contentplaceholder

\section{重参数化：如何穿过采样求梯度？}
\label{sec:vae-reparameterization}
% TODO：补充本节内容。
\contentplaceholder

\subsection{高斯重参数化与路径梯度}
\label{subsec:vae-pathwise}
% TODO：补充本节内容。
\contentplaceholder

\subsection{小批量训练流程}
\label{subsec:vae-training}
% TODO：补充本节内容。
\contentplaceholder

\subsection{先验采样与解码生成}
\label{subsec:vae-generation}
% TODO：补充本节内容。
\contentplaceholder

\section{为什么重建好还不够？}
\label{sec:vae-reconstruction}
% TODO：补充本节内容。
\contentplaceholder

\subsection{普通自编码器与潜在空间}
\label{subsec:vae-autoencoder}
% TODO：补充本节内容。
\contentplaceholder

\subsection{聚合后验与先验失配}
\label{subsec:vae-aggregated-posterior}
% TODO：补充本节内容。
\contentplaceholder

\subsection{重建与正则之间的取舍}
\label{subsec:vae-balance}
% TODO：补充本节内容。
\contentplaceholder

\section{重要变体与局限}
\label{sec:vae-variants}
% TODO：补充本节内容。
\contentplaceholder

\subsection{\texorpdfstring{$\beta$}{beta}-VAE：控制潜在信息容量}
\label{subsec:vae-beta-vae}
% TODO：补充本节内容。
\contentplaceholder

\subsection{后验坍塌：模型不再使用潜变量}
\label{subsec:vae-collapse}
% TODO：补充本节内容。
\contentplaceholder

\subsection{VQ-VAE：用码本学习离散表示}
\label{subsec:vae-vq-vae}
% TODO：补充本节内容。
\contentplaceholder

\subsection{码本损失、承诺损失与直通梯度}
\label{subsec:vae-vq-optimization}
% TODO：补充本节内容。
\contentplaceholder

\subsection{离散索引先验与 VQ-VAE-2}
\label{subsec:vae-discrete-prior}
% TODO：补充本节内容。
\contentplaceholder

\subsection{后续改进方向}
\label{subsec:vae-future}
% TODO：补充本节内容。
\contentplaceholder

\section{自检}
\label{sec:vae-review}
% TODO：补充本节内容。
\contentplaceholder

\section{来源与延伸阅读}
\label{sec:vae-reading}
% TODO：补充本节内容。
\contentplaceholder

% 来源：docs/06-normalizing-flows.md
% 仅建立章节框架；正文、公式、图表与习题内容待后续迁移。
\chapter{归一化流}
\label{chap:flow}

\section*{本章要点}
% TODO：迁移原章“先抓住这几件事”的要点。
\contentplaceholder

\section{为什么需要可逆生成模型？}
\label{sec:flow-motivation}
% TODO：补充本节内容。
\contentplaceholder

\subsection{确定性双射与同维潜变量}
\label{subsec:flow-bijection}
% TODO：补充本节内容。
\contentplaceholder

\subsection{基分布、生成方向与求逆方向}
\label{subsec:flow-directions}
% TODO：补充本节内容。
\contentplaceholder

\section{变量变换：保持概率质量，修正密度}
\label{sec:flow-change-of-variables}
% TODO：补充本节内容。
\contentplaceholder

\subsection{一维例子}
\label{subsec:flow-one-dimensional}
% TODO：补充本节内容。
\contentplaceholder

\subsection{多维公式}
\label{subsec:flow-multidimensional}
% TODO：补充本节内容。
\contentplaceholder

\subsection{局部体积变化与对数密度修正}
\label{subsec:flow-volume}
% TODO：补充本节内容。
\contentplaceholder

\section{多层组合与最大似然训练}
\label{sec:flow-composition}
% TODO：补充本节内容。
\contentplaceholder

\subsection{复合变换与对数行列式求和}
\label{subsec:flow-logdet}
% TODO：补充本节内容。
\contentplaceholder

\subsection{训练、生成与潜变量推断}
\label{subsec:flow-workflows}
% TODO：补充本节内容。
\contentplaceholder

\section{如何让 Jacobian 可计算？}
\label{sec:flow-jacobian}
% TODO：补充本节内容。
\contentplaceholder

\subsection{行列式计算成本与三角结构}
\label{subsec:flow-triangular}
% TODO：补充本节内容。
\contentplaceholder

\subsection{Planar Flow 与可逆性条件}
\label{subsec:flow-planar}
% TODO：补充本节内容。
\contentplaceholder

\subsection{NICE：加性耦合与缩放}
\label{subsec:flow-nice}
% TODO：补充本节内容。
\contentplaceholder

\subsection{Real NVP：仿射耦合}
\label{subsec:flow-real-nvp}
% TODO：补充本节内容。
\contentplaceholder

\section{自回归 Flow：MAF 与 IAF}
\label{sec:flow-autoregressive}
% TODO：补充本节内容。
\contentplaceholder

\subsection{MAF：条件参数依赖历史数据}
\label{subsec:flow-maf}
% TODO：补充本节内容。
\contentplaceholder

\subsection{IAF：条件参数依赖历史噪声}
\label{subsec:flow-iaf}
% TODO：补充本节内容。
\contentplaceholder

\subsection{生成、求逆与密度评估的效率比较}
\label{subsec:flow-comparison}
% TODO：补充本节内容。
\contentplaceholder

\section{概率密度蒸馏与 Parallel WaveNet}
\label{sec:flow-distillation}
% TODO：补充本节内容。
\contentplaceholder

\subsection{自回归教师与 IAF 学生}
\label{subsec:flow-teacher-student}
% TODO：补充本节内容。
\contentplaceholder

\subsection{反向 KL 蒸馏与计算流程}
\label{subsec:flow-reverse-kl}
% TODO：补充本节内容。
\contentplaceholder

\section{能力与限制}
\label{sec:flow-strengths-limits}
% TODO：补充本节内容。
\contentplaceholder

\subsection{精确密度、可逆结构与表达能力}
\label{subsec:flow-expressivity}
% TODO：补充本节内容。
\contentplaceholder

\subsection{离散数据与连续密度的区别}
\label{subsec:flow-discrete-data}
% TODO：补充本节内容。
\contentplaceholder

\subsection{连续时间 Flow 的后续阅读}
\label{subsec:flow-future}
% TODO：补充本节内容。
\contentplaceholder

\section{自检}
\label{sec:flow-review}
% TODO：补充本节内容。
\contentplaceholder

\section{来源与延伸阅读}
\label{sec:flow-reading}
% TODO：补充本节内容。
\contentplaceholder

% 来源：docs/07-generative-adversarial-networks.md
% 仅建立章节框架；正文、公式、图表与习题内容待后续迁移。
\chapter{生成对抗网络}
\label{chap:gan}

\section*{本章要点}
% TODO：迁移原章“先抓住这几件事”的要点。
\contentplaceholder

\section{为什么考虑不依赖似然的训练？}
\label{sec:gan-motivation}
% TODO：补充本节内容。
\contentplaceholder

\subsection{似然、视觉质量与样本多样性}
\label{subsec:gan-evaluation}
% TODO：补充本节内容。
\contentplaceholder

\subsection{双样本比较与分类器检验}
\label{subsec:gan-two-sample}
% TODO：补充本节内容。
\contentplaceholder

\section{两个网络与原始目标}
\label{sec:gan-objective}
% TODO：补充本节内容。
\contentplaceholder

\subsection{生成器与隐式分布}
\label{subsec:gan-generator}
% TODO：补充本节内容。
\contentplaceholder

\subsection{判别器与真假分类}
\label{subsec:gan-discriminator}
% TODO：补充本节内容。
\contentplaceholder

\subsection{原始 minimax 对抗目标}
\label{subsec:gan-minimax}
% TODO：补充本节内容。
\contentplaceholder

\section{最优判别器与 JS 散度}
\label{sec:gan-js}
% TODO：补充本节内容。
\contentplaceholder

\subsection{固定生成器时的最优判别器}
\label{subsec:gan-optimal-discriminator}
% TODO：补充本节内容。
\contentplaceholder

\subsection{JS 散度的定义与性质}
\label{subsec:gan-js-definition}
% TODO：补充本节内容。
\contentplaceholder

\subsection{最优判别器下目标的推导}
\label{subsec:gan-derivation}
% TODO：补充本节内容。
\contentplaceholder

\subsection{分布匹配的平衡点与理论条件}
\label{subsec:gan-equilibrium}
% TODO：补充本节内容。
\contentplaceholder

\section{实际训练：交替更新与非饱和损失}
\label{sec:gan-training}
% TODO：补充本节内容。
\contentplaceholder

\subsection{判别器与生成器的交替更新}
\label{subsec:gan-alternating}
% TODO：补充本节内容。
\contentplaceholder

\subsection{梯度饱和与非饱和生成器损失}
\label{subsec:gan-non-saturating}
% TODO：补充本节内容。
\contentplaceholder

\subsection{模式坍塌与训练不稳定}
\label{subsec:gan-collapse}
% TODO：补充本节内容。
\contentplaceholder

\section{WGAN：从分类差异到搬运距离}
\label{sec:gan-wgan}
% TODO：补充本节内容。
\contentplaceholder

\subsection{为什么 JS 可能缺乏距离信息？}
\label{subsec:gan-support}
% TODO：补充本节内容。
\contentplaceholder

\subsection{最优传输与对偶形式}
\label{subsec:gan-transport}
% TODO：补充本节内容。
\contentplaceholder

\subsection{Critic 的目标与 Lipschitz 约束}
\label{subsec:gan-critic}
% TODO：补充本节内容。
\contentplaceholder

\subsection{权重裁剪及其局限}
\label{subsec:gan-clipping}
% TODO：补充本节内容。
\contentplaceholder

\subsection{WGAN-GP 的梯度惩罚}
\label{subsec:gan-gradient-penalty}
% TODO：补充本节内容。
\contentplaceholder

\section{对抗损失可以与其他任务组合}
\label{sec:gan-applications}
% TODO：补充本节内容。
\contentplaceholder

\subsection{条件 GAN 与复合损失}
\label{subsec:gan-conditional}
% TODO：补充本节内容。
\contentplaceholder

\subsection{SRGAN、图像修复与 VQ-GAN}
\label{subsec:gan-restoration}
% TODO：补充本节内容。
\contentplaceholder

\subsection{Pix2Pix：成对图像翻译}
\label{subsec:gan-pix2pix}
% TODO：补充本节内容。
\contentplaceholder

\subsection{CycleGAN：非成对翻译与循环一致性}
\label{subsec:gan-cyclegan}
% TODO：补充本节内容。
\contentplaceholder

\subsection{StarGAN 与 ADDA：多域建模与域适应}
\label{subsec:gan-domains}
% TODO：补充本节内容。
\contentplaceholder

\section{从高分辨率生成到可控生成}
\label{sec:gan-controllability}
% TODO：补充本节内容。
\contentplaceholder

\subsection{Progressive Growing GAN}
\label{subsec:gan-progressive}
% TODO：补充本节内容。
\contentplaceholder

\subsection{StyleGAN 的中间潜在空间与样式控制}
\label{subsec:gan-stylegan}
% TODO：补充本节内容。
\contentplaceholder

\subsection{StyleGAN2 的结构改进}
\label{subsec:gan-stylegan2}
% TODO：补充本节内容。
\contentplaceholder

\subsection{StyleGAN3 与后续阅读}
\label{subsec:gan-future}
% TODO：补充本节内容。
\contentplaceholder

\section{自检}
\label{sec:gan-review}
% TODO：补充本节内容。
\contentplaceholder

\section{来源与延伸阅读}
\label{sec:gan-reading}
% TODO：补充本节内容。
\contentplaceholder

% 来源：docs/08-diffusion-models.md
% 仅建立章节框架；正文、公式、图表与习题内容待后续迁移。
\chapter{扩散模型：把生成拆成逐步去噪}
\label{chap:diffusion}

\section*{本章要点}
% TODO：迁移原章“先抓住这几件事”的要点。
\contentplaceholder

\section{前向加噪：一步就能构造任意噪声级别}
\label{sec:diffusion-forward}
% TODO：补充本节内容。
\contentplaceholder

\subsection{噪声日程与前向马尔可夫链}
\label{subsec:diffusion-schedule}
% TODO：补充本节内容。
\contentplaceholder

\subsection{任意时间步的闭式边缘分布}
\label{subsec:diffusion-marginal}
% TODO：补充本节内容。
\contentplaceholder

\subsection{末端高斯近似及其条件}
\label{subsec:diffusion-terminal}
% TODO：补充本节内容。
\contentplaceholder

\section{反向生成：从噪声到数据}
\label{sec:diffusion-reverse}
% TODO：补充本节内容。
\contentplaceholder

\subsection{反向链的联合分布}
\label{subsec:diffusion-joint}
% TODO：补充本节内容。
\contentplaceholder

\subsection{高斯反向转移与均值参数化}
\label{subsec:diffusion-gaussian}
% TODO：补充本节内容。
\contentplaceholder

\section{为什么训练时可以得到去噪监督？}
\label{sec:diffusion-supervision}
% TODO：补充本节内容。
\contentplaceholder

\subsection{三种条件分布的区别}
\label{subsec:diffusion-conditionals}
% TODO：补充本节内容。
\contentplaceholder

\subsection{前向后验的解析均值与方差}
\label{subsec:diffusion-posterior}
% TODO：补充本节内容。
\contentplaceholder

\subsection{训练监督与网络输入}
\label{subsec:diffusion-inputs}
% TODO：补充本节内容。
\contentplaceholder

\section{训练目标：一个多层潜变量模型的 ELBO}
\label{sec:diffusion-elbo}
% TODO：补充本节内容。
\contentplaceholder

\subsection{负对数似然的变分上界}
\label{subsec:diffusion-bound}
% TODO：补充本节内容。
\contentplaceholder

\subsection{终点先验、中间转移与重建项}
\label{subsec:diffusion-terms}
% TODO：补充本节内容。
\contentplaceholder

\subsection{从前向条件到反向 KL 的推导}
\label{subsec:diffusion-derivation}
% TODO：补充本节内容。
\contentplaceholder

\subsection{边界时间步与后验间隙}
\label{subsec:diffusion-endpoints}
% TODO：补充本节内容。
\contentplaceholder

\section{从预测均值到预测噪声}
\label{sec:diffusion-noise-prediction}
% TODO：补充本节内容。
\contentplaceholder

\subsection{噪声预测与反向均值的对应关系}
\label{subsec:diffusion-mean}
% TODO：补充本节内容。
\contentplaceholder

\subsection{高斯 KL 与加权噪声均方误差}
\label{subsec:diffusion-weighted-mse}
% TODO：补充本节内容。
\contentplaceholder

\subsection{简化噪声目标与时间步重加权}
\label{subsec:diffusion-simple-loss}
% TODO：补充本节内容。
\contentplaceholder

\section{训练与生成怎么进行？}
\label{sec:diffusion-algorithms}
% TODO：补充本节内容。
\contentplaceholder

\subsection{随机时间步的小批量训练}
\label{subsec:diffusion-training}
% TODO：补充本节内容。
\contentplaceholder

\subsection{时间条件与去噪网络架构}
\label{subsec:diffusion-network}
% TODO：补充本节内容。
\contentplaceholder

\subsection{逐步反向采样与最后一步处理}
\label{subsec:diffusion-sampling}
% TODO：补充本节内容。
\contentplaceholder

\subsection{干净数据估计与中间状态插值}
\label{subsec:diffusion-clean-estimate}
% TODO：补充本节内容。
\contentplaceholder

\section{与 VAE 的联系，以及方法的代价}
\label{sec:diffusion-connections}
% TODO：补充本节内容。
\contentplaceholder

\subsection{层次 VAE 视角与固定前向过程}
\label{subsec:diffusion-hierarchical-vae}
% TODO：补充本节内容。
\contentplaceholder

\subsection{样本质量、采样成本与数据类型}
\label{subsec:diffusion-cost}
% TODO：补充本节内容。
\contentplaceholder

\subsection{Score、SDE、DDIM、引导与潜在扩散的后续阅读}
\label{subsec:diffusion-future}
% TODO：补充本节内容。
\contentplaceholder

\section{自检}
\label{sec:diffusion-review}
% TODO：补充本节内容。
\contentplaceholder

\section{来源与延伸阅读}
\label{sec:diffusion-reading}
% TODO：补充本节内容。
\contentplaceholder


\appendix
% 来源：docs/formula-sheet.md
% 仅建立章节框架；正文、公式、图表与习题内容待后续迁移。
\chapter{核心公式速查}
\label{chap:formulas}

\section{概率分解与贝叶斯公式}
\label{sec:formulas-probability}
% TODO：补充本节内容。
\contentplaceholder

\section{最大似然与正向 KL}
\label{sec:formulas-mle}
% TODO：补充本节内容。
\contentplaceholder

\section{自回归似然}
\label{sec:formulas-autoregressive}
% TODO：补充本节内容。
\contentplaceholder

\section{ELBO 与后验间隙}
\label{sec:formulas-elbo}
% TODO：补充本节内容。
\contentplaceholder

\section{高斯 VAE 的实现}
\label{sec:formulas-gaussian-vae}
% TODO：补充本节内容。
\contentplaceholder

\section{Flow 的密度变换}
\label{sec:formulas-flow}
% TODO：补充本节内容。
\contentplaceholder

\section{GAN 与非饱和目标}
\label{sec:formulas-gan}
% TODO：补充本节内容。
\contentplaceholder

\section{Wasserstein 对偶}
\label{sec:formulas-wasserstein}
% TODO：补充本节内容。
\contentplaceholder

\section{扩散：前向加噪与噪声预测}
\label{sec:formulas-diffusion}
% TODO：补充本节内容。
\contentplaceholder

% 来源：docs/sources.md
% 仅建立章节框架；正文、公式、图表与习题内容待后续迁移。
\chapter{课件来源、覆盖范围与勘误}
\label{chap:sources}

\section{文件与知识章节的对应}
\label{sec:sources-mapping}
% TODO：补充本节内容。
\contentplaceholder

\section{如何取舍原始内容？}
\label{sec:sources-selection}
% TODO：补充本节内容。
\contentplaceholder

\section{公式勘误与严谨性说明}
\label{sec:sources-corrections}
% TODO：补充本节内容。
\contentplaceholder

\section{教材与课程入口}
\label{sec:sources-textbooks}
% TODO：补充本节内容。
\contentplaceholder

\section{按主题继续读哪些论文？}
\label{sec:sources-papers}
% TODO：补充本节内容。
\contentplaceholder

\subsection{图像统计与先验}
\label{subsec:sources-image-priors}
% TODO：补充本节内容。
\contentplaceholder

\subsection{自回归模型}
\label{subsec:sources-autoregressive}
% TODO：补充本节内容。
\contentplaceholder

\subsection{潜变量模型与 VAE}
\label{subsec:sources-vae}
% TODO：补充本节内容。
\contentplaceholder

\subsection{归一化流}
\label{subsec:sources-flow}
% TODO：补充本节内容。
\contentplaceholder

\subsection{生成对抗网络}
\label{subsec:sources-gan}
% TODO：补充本节内容。
\contentplaceholder

\subsection{扩散模型}
\label{subsec:sources-diffusion}
% TODO：补充本节内容。
\contentplaceholder


\end{document}
